\section{Integral lower degrees for the remaining alternating labels}
\label{sec:alternating-integral-cutoffs}

\paragraph{Status and scope.}
This is a new root corollary of the accepted exact shape inequalities
and all-count standard/two-box exclusions. The complete root proof/code
reading and the independent /root/sixfolds whole-proof audit have passed. It makes the existing necessary gate numerical; it is not
counted as an additional strict reduction beyond those exclusions.

In the main non-Eisenstein ordinary transitive alternating-core problem,
orient the partition, by transposing if necessary, as
$\lambda=(n-k,\mu)$ with $1\le k<n/4$ and $n\ge15$.
The accepted single-factor and all-count tensor exclusions now give
$k\ge3$. Put $h=2$ for the universal ALT07 bound. Where applicable,
one may instead use the larger values from ALT08: $h=4$ if eight
divides the original quadratic discriminant, or $h=p-1$ for an
odd discriminant prime $p\ge5$. Every such $h$ is an integer
at least two. The exact necessary inequality is
\[
 \beta(n,k)=\frac{k(n-1-k)}{(n-1)(n-2)}
       \le\frac1{5h}+\frac3{5e},
 \qquad e\ge\frac{n(n-3)}2.                         \tag{ALT9.1}
\]
Here $e$ is the actual inducing coefficient degree, not a chosen
upper bound; the degree floor is the reviewed nonstandard-label
bound used in ALT03/ALT08 and remains valid under tensor powers.

\begin{corollary}[HD-ALT-09]
\label{HD-ALT-09}
Under these hypotheses every surviving ordinary alternating label
satisfies the effective integral bound
\[
                              n\ge15h-1.                 \tag{ALT9.2}
\]
The integer is the smallest one allowed by the displayed numerical
conditions: at $n=15h-1$, the label $(n-3,3)$ satisfies
$\beta(n,3)<1/(5h)$. This last assertion is about necessary
inequalities only, not the existence of a curve source.
\end{corollary}

\begin{proof}
ALT07 already gives
$k<\tfrac12(1-\sqrt{3/5})n$. If $n\le26$, its right side is
at most $13(1-\sqrt{3/5})<3$: the strict last inequality follows
from $3\cdot169>5\cdot100$. Since $k\ge3$, necessarily
$n\ge27$.

On $3\le k<n/4$, the function $k(n-1-k)$ is increasing.
Consequently $\beta(n,k)\ge\beta(n,3)$.
Put $m=5h\ge10$ and suppose, for a contradiction, that
$27\le n\le3m-2$. Then $m\ge(n+2)/3$.
The desired strict inequality
\[
 \beta(n,3)-\frac1m>\frac6{5n(n-3)}                     \tag{ALT9.3}
\]
is equivalent, after multiplying positive denominators, to
\[
 F(n,m):=5n(n-3)\bigl(3m(n-4)-(n-1)(n-2)\bigr)
              -6m(n-1)(n-2)>0.                         \tag{ALT9.4}
\]
The coefficient of $m$ is
$15n(n-3)(n-4)-6(n-1)(n-2)>0$ for $n\ge27$.
For example, $(n-3)(n-4)>n$ in that range, so its first
term is greater than $15n^2$ and the subtracted term is
less than $6n^2$.
Thus it suffices to set $m=(n+2)/3$. Direct expansion gives
\[
 \begin{split}
 F\bigl(n,(n+2)/3\bigr)
  &=3n^3-63n^2+158n-8\\
  &=3(n-27)^3+180(n-27)^2+3317(n-27)+17380>0.
 \end{split}                                             \tag{ALT9.5}
\]
This proves (ALT9.3). The degree floor in (ALT9.1) gives
$3/(5e)\le6/(5n(n-3))$, so (ALT9.3) contradicts (ALT9.1).
Therefore $n\ge3m-1=15h-1$.

At the claimed numerical endpoint $n=3m-1$,
\[
 \beta(n,3)-\frac1m
             =\frac{-6}{m(n-1)(n-2)}<0.                 \tag{ALT9.6}
\]
By the same two-row hook formula used in ALT03,
the ordinary label $(n-3,3)$ has degree
$\binom n3-\binom n2=n(n-1)(n-5)/6$, at least
$n(n-3)/2$ since $(n-7)(n-2)\ge0$.
Also $3<n/4$ because $n\ge29$. It therefore satisfies
the numerical hypotheses and inequality at this endpoint,
including under any positive tensor power. This proves
the stated numerical minimality, without a source-existence claim.
\end{proof}

For the eight non-Eisenstein imaginary quadratic fields of class
number one, the resulting table is
\[
\begin{array}{c|r|r}
K&h&\text{necessary }n\text{ lower bound}\\\hline
\mathbf Q(i)&2&29\\
\mathbf Q(\sqrt{-2})&4&59\\
\mathbf Q(\sqrt{-7})&6&89\\
\mathbf Q(\sqrt{-11})&10&149\\
\mathbf Q(\sqrt{-19})&18&269\\
\mathbf Q(\sqrt{-43})&42&629\\
\mathbf Q(\sqrt{-67})&66&989\\
\mathbf Q(\sqrt{-163})&162&2429
\end{array}                                               \tag{ALT9.7}
\]
The exact tail and field inequalities remain necessary in addition
to these integers. This is an effective presentation of the remaining
condition-A local-system supplier gate (cote16, scan050;
SIGN_NEXT SN-C03/C04), with the same actual finite-cover,
representation lambda-ring and schemes-over-$\mathbf Z$ interpretation
as its antecedents. It gives no detecting cycle, does not invoke
condition-I positivity as a supplier, and adds no Markman, Severn
or032 premise.
